% layout: landscape
% TAB-30: generated by harness/analysis/paper_outputs.py; do not edit (every cell's source is in TAB-30.cells.json)
\begin{tabularx}{\linewidth}{l>{\hsize=0.967\hsize\raggedleft\hyphenpenalty=10000\exhyphenpenalty=10000\arraybackslash}X>{\hsize=0.901\hsize\raggedleft\hyphenpenalty=10000\exhyphenpenalty=10000\arraybackslash}X>{\hsize=1.081\hsize\raggedleft\hyphenpenalty=10000\exhyphenpenalty=10000\arraybackslash}X>{\hsize=0.455\hsize\raggedright\hyphenpenalty=10000\exhyphenpenalty=10000\arraybackslash}X>{\hsize=0.699\hsize\raggedleft\hyphenpenalty=10000\exhyphenpenalty=10000\arraybackslash}X>{\hsize=0.949\hsize\raggedleft\hyphenpenalty=10000\exhyphenpenalty=10000\arraybackslash}X>{\hsize=0.638\hsize\raggedleft\hyphenpenalty=10000\exhyphenpenalty=10000\arraybackslash}X>{\hsize=1.321\hsize\raggedleft\hyphenpenalty=10000\exhyphenpenalty=10000\arraybackslash}X>{\hsize=1.989\hsize\raggedright\hyphenpenalty=10000\exhyphenpenalty=10000\arraybackslash}X}
\toprule
Test & Price ratio per token (comparator / Haiku) & Samples per comparator answer & Cost-matched $k$ & Held & Vote accuracy at $k$ (\%) & Comparator accuracy (\%) & Gap $d$ (points) & 95\% interval & Source of $d$ and interval \\
\midrule
Test~1 & 10 & 11.01 & 11 &  & 64.44 & 100.00 & \ensuremath{-}35.56 & \mbox{[\ensuremath{-}44.85, \ensuremath{-}27.57]} & post~hoc (free~analysis~12) \\
Test~1 & 20 & 22.01 & 22 &  & 66.06 & 100.00 & \ensuremath{-}33.94 & \mbox{[\ensuremath{-}43.92, \ensuremath{-}24.85]} & post~hoc (free~analysis~12) \\
Test~1 & 40 & 44.03 & 40 & held at 40 & 66.88 & 100.00 & \ensuremath{-}33.13 & \mbox{[\ensuremath{-}43.75, \ensuremath{-}23.75]} & confirmatory (Test~1) \\
Test~1 & 80 & 88.05 & 40 & held at 40 & 66.88 & 100.00 & \ensuremath{-}33.13 & \mbox{[\ensuremath{-}43.75, \ensuremath{-}23.75]} & confirmatory (Test~1) \\
\addlinespace
Test~2 & 10 & 9.14 & 9 &  & 53.62 & 96.25 & \ensuremath{-}42.63 & \mbox{[\ensuremath{-}51.82, \ensuremath{-}33.00]} & post~hoc (free~analysis~12) \\
Test~2 & 20 & 18.29 & 18 &  & 56.71 & 96.25 & \ensuremath{-}39.54 & \mbox{[\ensuremath{-}49.82, \ensuremath{-}29.82]} & post~hoc (free~analysis~12) \\
Test~2 & 40 & 36.58 & 36 &  & 58.61 & 96.25 & \ensuremath{-}37.64 & \mbox{[\ensuremath{-}48.82, \ensuremath{-}26.96]} & confirmatory (Test~2) \\
Test~2 & 80 & 73.15 & 40 & held at 40 & 57.50 & 96.25 & \ensuremath{-}38.75 & \mbox{[\ensuremath{-}50.00, \ensuremath{-}28.13]} & post~hoc (free~analysis~12) \\
\addlinespace
Test~3 & 10 & 6.70 & 6 &  & 51.59 & 80.00 & \ensuremath{-}28.41 & \mbox{[\ensuremath{-}39.28, \ensuremath{-}16.13]} & post~hoc (free~analysis~12) \\
Test~3 & 20 & 13.40 & 13 &  & 55.30 & 80.00 & \ensuremath{-}24.70 & \mbox{[\ensuremath{-}36.12, \ensuremath{-}12.20]} & confirmatory (Test~3) \\
Test~3 & 40 & 26.80 & 26 &  & 58.21 & 80.00 & \ensuremath{-}21.79 & \mbox{[\ensuremath{-}34.31, \ensuremath{-}9.09]} & post~hoc (free~analysis~12) \\
Test~3 & 80 & 53.60 & 40 & held at 40 & 57.50 & 80.00 & \ensuremath{-}22.50 & \mbox{[\ensuremath{-}35.00, \ensuremath{-}8.75]} & post~hoc (free~analysis~12) \\
\addlinespace
Test~B2 & 10 & 1.68 & 1 &  & 92.63 & 97.81 & \ensuremath{-}5.19 & \mbox{[\ensuremath{-}7.75, \ensuremath{-}2.69]} & confirmatory (Test~B1: the same comparison as a one-sample vote) \\
Test~B2 & 20 & 3.35 & 3 &  & 97.03 & 97.81 & \ensuremath{-}0.79 & \mbox{[\ensuremath{-}2.83, 1.21]} & post~hoc (free~analysis~3 at this~$k$) \\
Test~B2 & 40 & 6.71 & 6 &  & 98.44 & 97.81 & +0.63 & \mbox{[\ensuremath{-}1.15, 2.44]} & confirmatory (Test~B2) \\
Test~B2 & 80 & 13.42 & 10 & held at 10 & 99.38 & 97.81 & +1.56 & \mbox{[\ensuremath{-}0.63, 3.44]} & post~hoc (free~analysis~3 at this~$k$) \\
\bottomrule
\end{tabularx}
% note: The stamped label of free analysis 12 reads ``post hoc sensitivity: how far the equal-money readings rest on one price list; no reading''; its ``equal-money'' is cost-matched.
% note: Samples per comparator answer: the comparator's cost per answer divided by the cost of one Haiku sample (at the price ratio in the row; in Panel~1, at list price).
% note: The price ratio is per token, comparator over Haiku, the same for input and output (the pinned table's is 40 for Opus and 20 for Sonnet); token counts are held as measured. At the list ratio each row is the locked test and prints its $d$ and interval.
% note: One comparison, one interval: any other re-priced row that is the same comparison as a printed result (the same settings, $k$ and puzzle set) takes that result and its label, named in the last column. A Test~1 row held at the locked test's $k$ is Test~1 as run; Test~B2's one-sample vote is Test~B1 (one Haiku-max sample against Opus-medium on the same 16-step set); its votes at other $k$ are free analysis 3's rows at that $k$. The remaining rows are post hoc (free analysis 12), with no reading.
% note: Held: the cost-matched $k$ at that price exceeds the samples drawn per puzzle and shape, so $k$ is held at that number.
% note: Accuracies in the main panel are given to two decimals, as the tests' own tables give them, so a row that is a locked test prints the same numbers; $d$ always has two. $d$ is the difference of the unrounded accuracies, so it can differ in its last digit from the difference of the two printed accuracies.
% note: Under the original pre-registered marking rule, which Deviation~2 kept as a sensitivity check, Test~3 would read inconclusive instead of worse by more than the margin ($d$~=~+0.30 points, 95\% interval \ensuremath{-}13.58 to +14.95).
% note: A break-even interval end is ``unbounded'' when too many replicates never reach the comparator's accuracy within the samples drawn; the share of such replicates is shown.
% note: Break-even intervals (Panel~1): the Haiku-max rows' cluster bootstrap uses 2,000 draws (free analysis 12); the Haiku-medium rows are Study~1's pre-specified break-even (item~9d), whose interval comes from the locked analysis, and the locked file does not stamp its number of draws, so none is given here.
% note: Cached view: the reading from the pre-registered decision rule, applied to the cached view's interval, is worse by more than the margin in every test (Tests~1--3).
% note: Cached view: Test~1's cached $k$ is held at the cap of 40 samples drawn; the cached cost of that many samples is still below the comparator's list cost per answer.
